\documentclass{article}

\usepackage{float}
\restylefloat{table}

\usepackage{booktabs}

\title{Team Productivity: POC\\\progname}

\author{\authname}

\date{}

\input{../Comments.text}
\input{../Common.text}

\begin{document}

\maketitle

This document summarizes the contributions of each team member up to the POC
Demo.  The time period of interest is the time between the beginning of the term
and the POC demo.

\noindent \textbf{Project Repository:}
\url{https://github.com/arjunkarthik123/fitness-finance-application}

\noindent \textbf{Data Collected:} November 17, 2026. The POC demos are
scheduled for the weeks of November 23 and November 30, 2026, so all counts
below cover September 18, 2026 (the first commit) to November 17, 2026.

\section{Demo Plans}

\subsection*{Feasibility Risk}

The main feasibility risk, as identified in the
\href{https://github.com/arjunkarthik123/fitness-finance-application/blob/main/pdfs/DevelopmentPlan/DevelopmentPlan.pdf}{Development
Plan}, is whether \progname{} can gamify and incentivize compounding habits by
rewarding healthy long-term thinking. This is the core thesis of the product.
The secondary risks are:
\begin{itemize}
    \item Improper workout verification, which leads to rewarding too little or
    too much virtual currency, or to users exploiting the tracking system.
    \item Users being unable to invest the virtual currency that they earn.
    \item Failure to integrate with physical fitness technologies.
\end{itemize}

\subsection*{What Will Be Demonstrated}

The POC will demonstrate the end-to-end flow that links fitness activity to
investing:
\begin{itemize}
    \item Connecting to a fitness data source (Android Health Connect) and
    importing the user's steps.
    \item Turning the imported fitness activity (steps) into virtual currency.
    \item Investing the virtual currency in a security.
    \item Simulating market behaviour for that security.
\end{itemize}

\section{Team Meeting Attendance}

\begin{table}[H]
\centering
\begin{tabular}{ll}
\toprule
\textbf{Student} & \textbf{Meetings}\\
\midrule
Total & 13\\
Arjun Karthik & 13\\
Shao (Daniel) Hu & 12\\
Noah Goldschmied & 12\\
Peter DeForest & 13\\
Daniel Franze-Da Silva & 13\\
\bottomrule
\end{tabular}
\end{table}

These counts cover the 13 team meetings held between September 27 and
November 14, 2026, which are recorded as \textit{Team Meeting} issues in the
repository.

\section{Supervisor/Stakeholder Meeting Attendance}

\noindent \textbf{Supervisor's Name: } Christopher Schankula

\begin{table}[H]
\centering
\begin{tabular}{ll}
\toprule
\textbf{Student} & \textbf{Meetings}\\
\midrule
Total & 2\\
Arjun Karthik & 2\\
Shao (Daniel) Hu & 2\\
Noah Goldschmied & 2\\
Peter DeForest & 2\\
Daniel Franze-Da Silva & 2\\
\bottomrule
\end{tabular}
\end{table}

The team's TA is also the team's supervisor (see the Team Meeting Plan in the
Development Plan). The team held 2 meetings with the supervisor, and every
member attended all of them. The user interviews used for the User Interviews
CD were held with potential users individually and are not counted here.

\section{Lecture Attendance}

\begin{table}[H]
\centering
\begin{tabular}{ll}
\toprule
\textbf{Student} & \textbf{Lectures}\\
\midrule
Total & 11\\
Arjun Karthik & 10\\
Shao (Daniel) Hu & 9\\
Noah Goldschmied & 10\\
Peter DeForest & 9\\
Daniel Franze-Da Silva & 11\\
\bottomrule
\end{tabular}
\end{table}

The total is the number of lectures held between the start of the term and
the POC demo.

\section{TA Document Discussion Attendance} \label{SecTA}

\noindent \textbf{TA's Name: } Christopher Schankula

\begin{table}[H]
\centering
\begin{tabular}{ll}
\toprule
\textbf{Student} & \textbf{Meetings}\\
\midrule
Total & 4\\
Arjun Karthik & 4\\
Shao (Daniel) Hu & 4\\
Noah Goldschmied & 4\\
Peter DeForest & 4\\
Daniel Franze-Da Silva & 4\\
\bottomrule
\end{tabular}
\end{table}

These counts come from the \textit{TA Meeting} issues for October 7, 2026 (SRS
and Hazard Analysis) and October 28, 2026 (VnV Plan).

\section{Commits, Lines Added and Lines Deleted, Pull Requests}

\begin{table}[H]
\centering
\begin{tabular}{lllllll}
\toprule
\textbf{Student} & \textbf{Commits} & \textbf{Percent} & \textbf{Lines Added} &
\textbf{Lines Deleted} & \textbf{Pull Requests}\\
\midrule
Total & 310 & 100\% & 22239 & 5303 & 171\\
Arjun Karthik & 44 & 14.2\% & 4281 & 821 & 34\\
Shao (Daniel) Hu & 58 & 18.7\% & 5799 & 2042 & 15\\
Noah Goldschmied & 49 & 15.8\% & 7785 & 1601 & 35\\
Peter DeForest & 79 & 25.5\% & 1752 & 561 & 49\\
Daniel Franze-Da Silva & 80 & 25.8\% & 2622 & 278 & 38\\
\midrule
GitHub Actions bot & 129 & -- & 21 & 20 & 0\\
\bottomrule
\end{tabular}
\end{table}

The counts were taken from the \texttt{main} branch with \texttt{git log} and
from the GitHub API, as follows:
\begin{itemize}
    \item \textbf{Commits} are non-merge commits on \texttt{main}. Merge commits
    created when pull requests were merged or branches were updated are not
    counted. The \texttt{github-actions} bot made 129 commits, which commit
    the compiled PDFs to the \texttt{pdfs/} folder. These are shown in a
    separate row and are not included in the team's total or percentages.
    One commit by the Copilot agent is also not counted. Commits made under more than one git identity (for example, a
    university email and a GitHub no-reply email) are combined for each
    member.
    \item \textbf{Lines Added and Lines Deleted} leave out generated and
    imported files: the initial commit that imported the course template
    (11,354 lines), \texttt{src/mobile/package-lock.json} (generated by npm,
    over 21,000 lines), PDF files, and the \texttt{.bst} bibliography style
    file. With these included, the totals are 67,055 lines added and 5,767
    lines deleted.
    \item \textbf{Pull Requests} are all pull requests opened by each member,
    including open pull requests and those closed without merging. Of the
    171 pull requests, 1 is open and 170 are closed, of which 166 were merged
    and 4 were closed without merging.
\end{itemize}

\section{Issue Tracker}

\begin{table}[H]
\centering
\begin{tabular}{lll}
\toprule
\textbf{Student} & \textbf{Authored (O+C)} & \textbf{Assigned (C only)}\\
\midrule
Arjun Karthik & 52 & 50 \\
Shao (Daniel) Hu & 44 & 57 \\
Noah Goldschmied & 40 & 53 \\
Peter DeForest & 48 & 46 \\
Daniel Franze-Da Silva & 51 & 57 \\
\bottomrule
\end{tabular}
\end{table}

These counts include issues only, not pull requests. 18 of the closed issues
(the meeting, lecture and TA meeting issues) are assigned to all five members,
so they are included in every member's assigned count. The 19 peer review
issues opened by Team 9 are not included in the authored counts, since they
were not written by members of this team.

\section{CICD via GitHub Actions}

The team uses two GitHub Actions workflows:
\begin{itemize}
    \item \textbf{App CI}
    (\href{https://github.com/arjunkarthik123/fitness-finance-application/blob/main/.github/workflows/app-ci.yml}{\texttt{app-ci.yml}}):
    runs on every push and pull request to \texttt{main} that changes the
    source code. The \texttt{mobile} job installs the React Native app's
    dependencies and runs ESLint, the TypeScript type check, the Prettier
    format check and the Jest tests. The \texttt{api} job starts a MySQL
    database and runs \texttt{./mvnw verify} on the Spring Boot backend, which
    builds the API and runs its tests against the same database used in
    production.
    \item \textbf{Build LaTeX and Deploy PDFs}
    (\href{https://github.com/arjunkarthik123/fitness-finance-application/blob/main/.github/workflows/latex-pages.yml}{\texttt{latex-pages.yml}}):
    compiles every changed LaTeX document on each pull request. The
    \texttt{build} check is required before a pull request can be merged, so
    a document that does not compile cannot reach \texttt{main}. On a push to
    \texttt{main}, the workflow commits the compiled PDFs to the
    \texttt{pdfs/} folder for TA review and deploys them to GitHub Pages.
\end{itemize}

As development continues, the team plans to add the remaining unit and
integration tests to these workflows as each module is implemented, and to
add the end-to-end tests once the main screens are in place.

\section{Custom Documents (CDs)}


The team has completed the User Interviews CD
(\href{https://github.com/arjunkarthik123/fitness-finance-application/blob/main/docs/CDs/Interviews/UserInterviewsFall.tex}{\texttt{UserInterviewsFall.tex}}),
which summarizes the interviews with potential users, includes an empathy map
of a novice participant, and lists the key takeaways that informed the SRS.

\section{Team Charter Trigger Items}

\subsection*{Quantified Triggers}

The team charter (Development Plan, Appendix --- Team Charter) defines the
following quantified triggers:
\begin{itemize}
    \item Each member attends at least 90\% of all scheduled meetings.
    \item Missing 2 meetings in a row without an acceptable excuse or proper
    communication leads to an emergency team meeting, and to escalation to the
    TA and/or instructor if the issue continues.
    \item Arriving late or leaving early without an acceptable excuse more
    than 3 times leads to an emergency team meeting, and to escalation if the
    issue continues.
    \item Members notify the team at least 24 hours before missing a meeting,
    unless it is an emergency.
    \item By the end of the project, each member contributes 20--30\% of the
    total contributions.
\end{itemize}

\subsection*{Violations}

There have been no violations of the triggers. Across the 30 scheduled
meetings (13 team meetings, 11 lectures, 4 TA meetings and 2 supervisor
meetings), every member attended at least 90\%, which meets the attendance
trigger.

\subsection*{Plan to Address the Violations}

As there are no violations, no corrective action is needed. To keep the
triggers easy to check, the team will add a field to the team meeting and
lecture issue templates for the reason for an absence, and will review
attendance again at the end of the fall term.

\section{Additional Productivity Metrics}

\begin{table}[H]
\centering
\begin{tabular}{ll}
\toprule
\textbf{Metric} & \textbf{Value}\\
\midrule
Pull requests open & 1 of 171\\
Pull requests closed & 170 of 171 (166 merged)\\
Issues closed & 210 of 254\\
Peer review issues received (Team 9) & 19 (0 closed)\\
TA feedback issues created & 19 (0 closed)\\
Automated test files (mobile app) & 6\\
Automated test files (API) & 3\\
\bottomrule
\end{tabular}
\end{table}

The TA feedback issues were created on November 14--15, 2026, and will be
addressed in Revision 1 of each document, along with the peer review issues.

\subsection*{Pull Request Reviews}

\begin{table}[H]
\centering
\begin{tabular}{ll}
\toprule
\textbf{Student} & \textbf{PRs Reviewed}\\
\midrule
Arjun Karthik & 39\\
Shao (Daniel) Hu & 13\\
Noah Goldschmied & 22\\
Peter DeForest & 44\\
Daniel Franze-Da Silva & 11\\
\bottomrule
\end{tabular}
\end{table}

These counts are the number of pull requests on which each member submitted at
least one review (an approval, a comment or a request for changes), taken from
GitHub. Each pull request is counted once per reviewer. Comments that were not
part of a review, and reviews by Copilot or other bots, are not counted.

Not every pull request was reviewed. As recorded in the October 18, 2026 team
meeting issue, the team temporarily disabled branch protection and mandatory
reviews while finishing the SRS and Hazard Analysis, because of other school
deadlines and a busy schedule for all team members. Instead of reviewing each
pull request individually, the team reviewed the SRS and every member's
sections together during that meeting.

\end{document}
